\documentclass[14pt,a4paper,oneside]{extbook}
\usepackage[utf8]{vietnam}
\usepackage{amsmath,amssymb,mathrsfs,fancyhdr,enumerate,multirow,makecell}
\usepackage{tkz-euclide,tikz-3dplot,tikz,tkz-tab}
\usetikzlibrary{arrows.meta,calc,intersections,angles,quotes,backgrounds,shapes.geometric,patterns,through,shadows,positioning}
\usepackage{yhmath}
\usepackage[top=0.7cm, bottom=0.7cm, left=1.0cm, right=1.0cm, headheight=16pt, headsep=3mm, footskip=5mm]{geometry}
\usepackage[hidelinks,unicode]{hyperref}
\usepackage{qrcode}
\usepackage[many]{tcolorbox}
\usepackage{fontawesome}

\pagestyle{fancy}
\fancyhf{}
\fancyhead[L]{\footnotesize\bfseries\color{teal!85!black}\faBook~SỔ TAY KIẾN THỨC TOÁN 5}
\fancyhead[R]{\footnotesize\bfseries\color{blue!85!black}Lớp Toán Thầy Lĩnh -- loptoanthaylinh.vn}
\fancyfoot[C]{\footnotesize\bfseries\color{gray!80!black}-- Sổ tay kiến thức: Bài 3. Ôn tập phân số --}
\renewcommand{\headrulewidth}{0.4pt}

\begin{document}

% --- BANNER TIÊU ĐỀ BÀI HỌC ---
\noindent
\begin{tcolorbox}[
  enhanced,
  colback=teal!8!white,
  colframe=teal!80!black,
  arc=3mm,
  boxrule=1.1pt,
  width=\linewidth,
  halign=center,
  top=2.2mm, bottom=2.2mm,
  left=3mm, right=3mm,
  drop fuzzy shadow
]
  {\large\bfseries\color{teal!95!black}\faBookmark\ BÀI 03. ÔN TẬP PHÂN SỐ}\\[2pt]
  {\scriptsize\color{black!75}\itshape Khái niệm phân số, rút gọn, quy đồng mẫu số và so sánh phân số trên 1 trang A4}
\end{tcolorbox}

\vspace{3.2mm}
% --- KHỐI 1: KHÁI NIỆM & TÍNH CHẤT CƠ BẢN ---
\noindent
\begin{minipage}[t]{0.56\textwidth}
  \begin{tcolorbox}[
    enhanced,
    title={\footnotesize\bfseries\color{white}\faKey\ 1. Khái niệm và tính chất cơ bản},
    colback=blue!2!white,
    colframe=blue!75!black,
    arc=2.5mm, boxrule=0.8pt, top=1.8mm, bottom=1.8mm, left=2.2mm, right=2.2mm
  ]
    \footnotesize
    $\bullet$ Phân số có dạng $\dfrac{a}{b}$ ($a$: tử số, $b$: mẫu số, $b \neq 0$).\\
    $\bullet$ \textbf{Tính chất cơ bản:}\\
    $\triangleright$ Nhân cả tử và mẫu với cùng số $m \neq 0$: $\dfrac{a}{b} = \dfrac{a \times m}{b \times m}$.\\
    $\triangleright$ Chia cả tử và mẫu cho ước chung $n$: $\dfrac{a}{b} = \dfrac{a : n}{b : n}$.\\
    $\bullet$ Phân số tối giản: Tử số và mẫu số không cùng chia hết cho số tự nhiên nào lớn hơn $1$.
  \end{tcolorbox}
\end{minipage}%
\hfill
\begin{minipage}[t]{0.42\textwidth}
  \begin{tcolorbox}[
    enhanced,
    title={\footnotesize\bfseries\color{white}\faCrop\ 2. Rút gọn \& Quy đồng},
    colback=teal!2!white,
    colframe=teal!75!black,
    arc=2.5mm, boxrule=0.8pt, top=1.8mm, bottom=1.8mm, left=2mm, right=2mm
  ]
    \footnotesize
    $\bullet$ \textbf{Rút gọn:} Chia cả tử và mẫu cho một số lớn hơn $1$ để được phân số tối giản.\\
    \textit{VD:} $\dfrac{18}{24} = \dfrac{18:6}{24:6} = \dfrac{3}{4}$.\\[2pt]
    $\bullet$ \textbf{Quy đồng mẫu số:}\\
    $\triangleright$ Chọn mẫu số chung (MSC).\\
    $\triangleright$ Nhân cả tử và mẫu với thừa số phụ tương ứng.
  \end{tcolorbox}
\end{minipage}

\vspace{3mm}

% --- KHỐI 2: CÁC CÁCH SO SÁNH PHÂN SỐ ---
\noindent
\begin{tcolorbox}[
  enhanced,
  title={\footnotesize\bfseries\color{white}\faBalanceScale\ 3. Bốn phương pháp so sánh phân số},
  colback=purple!2!white,
  colframe=purple!70!black,
  arc=2.5mm, boxrule=0.8pt, top=1.8mm, bottom=1.8mm, left=2.5mm, right=2.5mm
]
  \footnotesize
  \begin{tabular}{ll}
    $\bullet$ \textbf{Cùng mẫu số ($b>0$):} Phân số nào có tử số lớn hơn thì lớn hơn: $\dfrac{3}{7} < \dfrac{5}{7}$.\\
    $\bullet$ \textbf{Cùng tử số ($a>0$):} Phân số nào có mẫu số bé hơn thì lớn hơn: $\dfrac{4}{5} > \dfrac{4}{9}$.\\
    $\bullet$ \textbf{Khác cả tử và mẫu:} Quy đồng mẫu số (hoặc quy đồng tử số) rồi so sánh.\\
    $\bullet$ \textbf{So sánh qua số trung gian (thường là $1$):}\\
    \quad $\triangleright$ Phân số bé hơn $1$ (tử $<$ mẫu) so với phân số lớn hơn $1$ (tử $>$ mẫu): $\dfrac{7}{9} < 1 < \dfrac{11}{8} \Rightarrow \dfrac{7}{9} < \dfrac{11}{8}$.\\
    \quad $\triangleright$ So sánh qua phần bù đơn vị ($1 - \dfrac{a}{b}$) hoặc phần hơn đơn vị ($\dfrac{a}{b} - 1$).
  \end{tabular}
\end{tcolorbox}

\vspace{3mm}

% --- KHỐI 3: SƠ ĐỒ TƯ DUY ---
\noindent
\begin{tcolorbox}[
  enhanced,
  title={\footnotesize\bfseries\color{white}\faSitemap\ Sơ đồ tư duy: Kĩ năng xử lý phân số},
  colback=yellow!2!white,
  colframe=teal!85!black,
  arc=3mm, boxrule=0.9pt, top=1.8mm, bottom=1.8mm, left=1.5mm, right=1.5mm,
  halign=center
]
  \begin{tikzpicture}[
    scale=0.92,
    every node/.style={transform shape},
    root/.style={rectangle, rounded corners=5pt, fill=teal!85!black, text=white, font=\bfseries\footnotesize, inner sep=5pt},
    b1/.style={rectangle, rounded corners=4pt, fill=blue!12, draw=blue!75!black, text=blue!90!black, font=\bfseries\scriptsize, inner sep=4pt},
    b2/.style={rectangle, rounded corners=4pt, fill=purple!12, draw=purple!75!black, text=purple!90!black, font=\bfseries\scriptsize, inner sep=4pt},
    b3/.style={rectangle, rounded corners=4pt, fill=orange!12, draw=orange!75!black, text=orange!90!black, font=\bfseries\scriptsize, inner sep=4pt},
    b4/.style={rectangle, rounded corners=4pt, fill=green!15, draw=green!65!black, text=green!80!black, font=\bfseries\scriptsize, inner sep=4pt}
  ]
    \node[root] (core) at (0,0) {\faBook\ PHÂN SỐ};
    \node[b1] (n1) at (-4.8, 0.8) {\faSearch\ Rút gọn tối giản};
    \node[b2] (n2) at (4.8, 0.8) {\faThList\ Quy đồng mẫu số};
    \node[b3] (n3) at (-4.8, -0.8) {\faBalanceScale\ So sánh phân số};
    \node[b4] (n4) at (4.8, -0.8) {\faSliders\ Phần bù / Phần hơn};

    \draw[teal!85!black, thick, -{Stealth[length=2mm]}] (core) -- (n1);
    \draw[teal!85!black, thick, -{Stealth[length=2mm]}] (core) -- (n2);
    \draw[teal!85!black, thick, -{Stealth[length=2mm]}] (core) -- (n3);
    \draw[teal!85!black, thick, -{Stealth[length=2mm]}] (core) -- (n4);
  \end{tikzpicture}
\end{tcolorbox}

\vspace{3mm}

% --- KHỐI 4: MẸO SO SÁNH NHANH & BẪY TRẮC NGHIỆM ---
\noindent
\begin{minipage}[t]{0.54\textwidth}
  \begin{tcolorbox}[
    enhanced,
    title={\footnotesize\bfseries\color{white}\faBolt\ Mẹo nhân chéo so sánh siêu tốc},
    colback=green!3!white,
    colframe=green!60!black,
    arc=2.5mm, boxrule=0.8pt, top=1.8mm, bottom=1.8mm, left=2.2mm, right=2.2mm
  ]
    \footnotesize
    Để so sánh hai phân số dương $\dfrac{a}{b}$ và $\dfrac{c}{d}$:\\
    $\bullet$ Tính tích chéo: $a \times d$ và $b \times c$.\\
    $\triangleright$ Nếu $a \times d > b \times c$ thì $\dfrac{a}{b} > \dfrac{c}{d}$.\\
    $\triangleright$ Nếu $a \times d < b \times c$ thì $\dfrac{a}{b} < \dfrac{c}{d}$.\\
    \textit{VD:} So sánh $\dfrac{4}{7}$ và $\dfrac{5}{9}$: Ta có $4 \times 9 = 36 > 7 \times 5 = 35 \Rightarrow \dfrac{4}{7} > \dfrac{5}{9}$.
  \end{tcolorbox}
\end{minipage}%
\hfill
\begin{minipage}[t]{0.44\textwidth}
  \begin{tcolorbox}[
    enhanced,
    title={\footnotesize\bfseries\color{white}\faExclamationTriangle\ Mẹo tránh bẫy trắc nghiệm},
    colback=red!2!white,
    colframe=red!70!black,
    arc=2.5mm, boxrule=0.8pt, top=1.8mm, bottom=1.8mm, left=2.2mm, right=2.2mm
  ]
    \footnotesize
    $\bullet$ Rất nhiều học sinh nhầm: Cùng tử số thì mẫu lớn hơn lại tưởng là lớn hơn! \textbf{Ghi nhớ}: Mẫu càng lớn, mỗi phần chia càng nhỏ.\\
    $\bullet$ Mọi số tự nhiên đều viết được dưới dạng phân số có mẫu số bằng $1$: $n = \dfrac{n}{1}$.
  \end{tcolorbox}
\end{minipage}

\end{document}
